\documentclass[12pt,a4paper]{article}
\usepackage[utf8]{inputenc}
\usepackage[spanish]{babel}
\usepackage{graphicx}
\usepackage{listings}
\usepackage[table]{xcolor}
\usepackage{makecell}
\usepackage{tikz}
\usepackage{booktabs}
\usepackage{geometry}
\usepackage{hyperref}
\usepackage{fancyhdr}
\usepackage{titlesec}

\geometry{margin=2.5cm}
\hypersetup{colorlinks=true,linkcolor=blue,urlcolor=blue}
\setlength{\parskip}{0.3cm}

\definecolor{codebg}{RGB}{240,240,240}
\definecolor{codekey}{RGB}{0,0,200}
\definecolor{codestr}{RGB}{0,130,0}
\definecolor{codecmt}{RGB}{140,140,140}

\lstset{
  language=Python,
  basicstyle=\ttfamily\footnotesize,
  backgroundcolor=\color{codebg},
  keywordstyle=\color{codekey}\bfseries,
  stringstyle=\color{codestr},
  commentstyle=\color{codecmt}\itshape,
  showstringspaces=false,
  frame=single,
  framerule=0.5pt,
  aboveskip=6pt,
  belowskip=6pt,
  tabsize=2,
  breaklines=true,
}

\pagestyle{fancy}
\fancyhf{}
\fancyhead[L]{Guía de Implementación -- Juego de la Vida}
\fancyhead[R]{Equipo Rosa -- PPyC 2026}
\fancyfoot[C]{\thepage}

\begin{document}

% ================================================================
% PORTADA
% ================================================================
\begin{titlepage}
  \centering
  \vspace*{3cm}

  {\Huge \textbf{Guía de Implementación}}

  \vspace{0.5cm}
  {\Large Juego de la Vida de Conway}

  \vspace{0.3cm}
  {\large Arquitectura en 3 capas}

  \vspace{1.5cm}

  \begin{tabular}{rl}
    \textbf{Capa Fundación}      & Israel Olvera Juárez \\
    \textbf{Capa Refinamiento}   & Miguel Ángel Torres Pérez \\
    \textbf{Capa Escalabilidad}  & Andrés Soria Cabrera \\
  \end{tabular}

  \vspace{2cm}

  {\large Programación Paralela y Concurrente}

  {\large Universidad Nacional Autónoma de México}

  \vspace{0.3cm}
  {\large Mayo 2026}

  \vspace{0.5cm}

  \begin{center}
  \renewcommand{\arraystretch}{2.5}
  \begin{tabular}{|c|}
    \hline
    \cellcolor{green!15}
    \textbf{\large Fundación}  \\
    \cellcolor{green!15}
    Israel Olvera Juárez  \\
    \cellcolor{green!15}
    \texttt{juego\_vida.py} \\
    \hline
    \multicolumn{1}{c}{$\big\downarrow$} \\
    \hline
    \cellcolor{yellow!15}
    \textbf{\large Refinamiento} \\
    \cellcolor{yellow!15}
    Miguel Ángel Torres Pérez \\
    \cellcolor{yellow!15}
    \texttt{juego\_vida\_optimo.py} \\
    \hline
    \multicolumn{1}{c}{$\big\downarrow$} \\
    \hline
    \cellcolor{red!15}
    \textbf{\large Escalabilidad} \\
    \cellcolor{red!15}
    Andrés Soria Cabrera \\
    \cellcolor{red!15}
    \texttt{juego\_vida\_proyecto.py} \\
    \hline
  \end{tabular}
  \end{center}
\end{titlepage}

\newpage

% ================================================================
% SECCIÓN 0: VISIÓN GENERAL
% ================================================================
\section{Visión general de la arquitectura}

\subsection{Las 3 capas}

El proyecto está organizado en tres capas que se construyen una sobre otra:

\begin{center}
\renewcommand{\arraystretch}{2.2}
\begin{tabular}{ccl}
  \cellcolor{green!15}
  \begin{minipage}{4cm}\centering
    \textbf{Capa 1: Fundación} \\
    \texttt{juego\_vida.py}
  \end{minipage}
  &
  $\rightarrow$
  &
  \begin{minipage}{5cm}\raggedright\footnotesize
    GUI Pygame, Set[Cell], \\
    vecindario de Moore, \\
    algoritmo de 2 pasadas
  \end{minipage}
  \\
  \cellcolor{yellow!15}
  \begin{minipage}{4cm}\centering
    \textbf{Capa 2: Refinamiento} \\
    \texttt{juego\_vida\_optimo.py}
  \end{minipage}
  &
  $\rightarrow$
  &
  \begin{minipage}{5cm}\raggedright\footnotesize
    Métodos puros, API programable, \\
    catálogo de 21 patrones, \\
    serialización YAML
  \end{minipage}
  \\
  \cellcolor{red!15}
  \begin{minipage}{4cm}\centering
    \textbf{Capa 3: Escalabilidad} \\
    \texttt{juego\_vida\_proyecto.py}
  \end{minipage}
  &
  $\rightarrow$
  &
  \begin{minipage}{5cm}\raggedright\footnotesize
    Multiprocessing, Pool.map, \\
    11 reglas B/S configurables, \\
    cámara y renderizado
  \end{minipage}
  \\
\end{tabular}
\end{center}

\subsection{Archivos del proyecto}

\begin{table}[h]
\centering
\begin{tabular}{llp{6cm}}
\toprule
\textbf{Archivo} & \textbf{Tipo} & \textbf{Propósito} \\
\midrule
\texttt{config.yml}       & Compartido & Dimensiones de grilla, colores, regla por defecto \\
\texttt{requirements.txt} & Compartido & Dependencias (pygame, numpy, PyYAML, scipy, etc.) \\
\texttt{seed\_patterns.yml} & Fundación & 9 patrones en matriz binaria (teclas 1--9) \\
\texttt{juego\_vida.py}     & Fundación & Implementación base con GUI Pygame \\
\texttt{patterns.py}        & Refinamiento & Catálogo de 21 patrones + funciones de serialización \\
\texttt{juego\_vida\_optimo.py} & Refinamiento & Núcleo algorítmico puro, sin GUI \\
\texttt{automata.py}        & Escalabilidad & Motor de reglas B/S con 11 presets \\
\texttt{juego\_vida\_proyecto.py} & Escalabilidad & Versión paralela con multiprocessing \\
\bottomrule
\end{tabular}
\end{table}

\subsection{Configuración central: \texttt{config.yml}}

Todas las implementaciones leen este archivo. Los campos principales son:

\begin{table}[h]
\centering
\footnotesize
\begin{tabular}{p{4.5cm}p{2cm}p{7cm}}
\toprule
\textbf{Campo} & \textbf{Valor} & \textbf{Descripción} \\
\midrule
\texttt{width / height} & 800 & Tamaño de ventana en píxeles \\
\texttt{cell\_size}      & 16  & Píxeles por célula ($50\times50$ celdas) \\
\texttt{grid\_is\_infinite} & True & Wrapping toroidal con operador \% \\
\texttt{rule\_string}    & B3/S23 & Regla de autómata por defecto \\
\texttt{sleep}           & 0.15 & Segundos entre generaciones (solo v1) \\
\texttt{fps}             & 40   & Fotogramas por segundo \\
\makecell[l]{\texttt{reproduction}\\ \texttt{underpopulation}\\ \texttt{overpopulation}} & 3/2/3 & B3/S23 clásico (legado) \\
\bottomrule
\end{tabular}
\end{table}

\subsection{Dependencias entre módulos}

\begin{verbatim}
config.yml  <── leído por los 3 .py
     │
     ├── automata.py ──── (sin dependencias del proyecto)
     │      │
     └── patterns.py ──── (sin dependencias del proyecto)
            │
     ┌──────┴───────────────┬──────────────────────┐
     │                      │                      │
juego_vida.py        juego_vida_optimo.py    juego_vida_proyecto.py
(GUI Pygame)         (headless)              (multiprocessing)
\end{verbatim}

\subsection{Tipo común}

El tipo \texttt{Cell = Tuple[int, int]} representa una coordenada \texttt{(columna, fila)}
en la grilla. Aparece en \texttt{patterns.py} y es usado por todas las implementaciones.

\newpage

% ================================================================
% SECCIÓN 1: FUNDACIÓN
% ================================================================
\section{Capa Fundación -- Israel Olvera Juárez}

\subsection{Archivos asignados}

\begin{itemize}
  \item \texttt{juego\_vida.py} -- implementación base con GUI Pygame (259 líneas)
  \item \texttt{seed\_patterns.yml} -- 9 patrones en formato de matriz binaria
\end{itemize}

\subsection{Clase GameOfLife}

La clase \texttt{GameOfLife} contiene toda la lógica del juego y el renderizado.

\subsubsection{Constructor}

\begin{lstlisting}
def __init__(self) -> None:
    self._config = self._get_config()      # lee config.yml
    self._n_cols = 800 // 16               # = 50 columnas
    self._n_rows = 800 // 16               # = 50 filas
    self._living_cells: Set[Cell] = set()  # conjunto disperso
    self._is_paused = True                 # empieza en pausa
    self._rule = get_rule("B3/S23")        # regla de Conway
\end{lstlisting}

\subsubsection{La decisión de diseño más importante: \texttt{Set[Cell]}}

En vez de una matriz 2D de 50$\times$50 = 2,500 celdas, solo se almacenan
las coordenadas de las células \textbf{vivas}. Esto hace que:

\begin{itemize}
  \item El paso de simulación sea $O(\text{vivas} + \text{vecinas})$ en vez de $O(50 \times 50)$
  \item En grillas dispersas (pocas células) sea extremadamente rápido
  \item La memoria usada sea mínima
\end{itemize}

Una célula viva en la columna 5, fila 10 es simplemente la tupla \texttt{(5, 10)}.
Una célula muerta \textbf{no se almacena}; se asume muerta si no está en el conjunto.

\subsubsection{Algoritmo de 2 pasadas -- \texttt{run\_logic()}}

Este es el corazón del juego. Cada generación se calcula en dos pasos:

\begin{enumerate}
  \item \textbf{Paso 1 -- Supervivencia}: para cada célula viva, contar sus 8 vecinos.
        Si \texttt{should\_survive(vecinos)} es True, se mantiene viva.
        Además, se recolectan todos los vecinos de células vivas en \texttt{all\_neighbors}.

  \item \textbf{Paso 2 -- Reproducción}: para cada vecino de una célula viva que esté
        \textbf{muerto}, contar sus vecinos. Si \texttt{should\_reproduce(vecinos)} es True, nace.
\end{enumerate}

\begin{lstlisting}
def run_logic(self) -> None:
    new_living_cells = set()
    all_neighbors = set()

    # Paso 1: supervivencia
    for cell in self._living_cells:
        n = self._count_living_neighbors(cell)
        if self._should_cell_survive(n):
            new_living_cells.add(cell)
        all_neighbors.update(self._get_neighbors(cell))

    # Paso 2: reproducción
    for cell in all_neighbors:
        if cell not in self._living_cells:
            n = self._count_living_neighbors(cell)
            if self._should_cell_reproduce(n):
                new_living_cells.add(cell)

    self._living_cells = new_living_cells
\end{lstlisting}

\subsubsection{Vecindario de Moore con wrapping}

Cada célula tiene 8 vecinos (horizontal, vertical y diagonal). Con \texttt{grid\_is\_infinite = True},
los bordes se envuelven usando el operador módulo (\texttt{\%}), creando una grilla toroidal.

\begin{lstlisting}
def _get_neighbors(self, cell: Cell) -> List[Cell]:
    col, row = cell
    neighbors = []
    for dc in [-1, 0, 1]:
        for dr in [-1, 0, 1]:
            if dc == 0 and dr == 0:
                continue
            nc = (col + dc) % self._n_cols    # wrapping
            nr = (row + dr) % self._n_rows
            neighbors.append((nc, nr))
    return neighbors
\end{lstlisting}

\subsubsection{Controles de usuario}

\begin{table}[h]
\centering
\begin{tabular}{lll}
\toprule
\textbf{Tecla / Acción} & \textbf{Método} & \textbf{Efecto} \\
\midrule
Espacio        & --                 & Pausar / reanudar simulación \\
Flecha derecha & --                 & Avanzar 1 solo paso (en pausa) \\
C              & \texttt{clear()}   & Limpiar todas las células \\
G / 0          & \texttt{\_generate\_random\_init\_grid()} & Grilla aleatoria (5\%--20\% densidad) \\
1--9           & \texttt{\_generate\_pattern\_from\_catalog()} & Cargar patrón del catálogo \\
P              & \texttt{\_print\_catalog()} & Imprimir catálogo en consola \\
R              & \texttt{\_cycle\_rule()} & Cambiar regla de autómata \\
S              & \texttt{\_save\_current\_pattern()} & Guardar patrón a YAML \\
Click mouse    & --                 & Alternar célula viva/muerta \\
\bottomrule
\end{tabular}
\end{table}

\subsection{\texttt{seed\_patterns.yml}}

Contiene 9 patrones (blinker, glider, lwss, pulsar, gosper\_glider\_gun, toad, beacon,
r\_pentomino, diehard) en formato de matriz binaria. Cada patrón tiene un ID numérico
que corresponde a la tecla que lo carga.

Ejemplo del formato:
\begin{lstlisting}
- patterns:
    1:
      name: "blinker"
      matrix: [[0,1,0], [0,1,0], [0,1,0]]
\end{lstlisting}

La función \texttt{\_generate\_pattern\_from\_catalog(id\_)} busca el nombre del patrón
en un diccionario (\texttt{\_PATTERN\_KEY\_MAP}) y lo carga del catálogo de \texttt{patterns.py}.

\subsection{Puntos para exponer}

\begin{enumerate}
  \item \textbf{La representación Set[Cell]}: explicar por qué es más eficiente que una matriz.
        Solo guardamos lo que está vivo. Una grilla de 50$\times$50 con 10\% de densidad
        almacena solo 250 tuplas en vez de 2,500 enteros.
  \item \textbf{El algoritmo de 2 pasadas}: primero decidimos qué células vivas sobreviven,
        luego qué células muertas nacen. Es simple pero poderoso.
  \item \textbf{El vecindario de Moore}: 8 vecinos, y con wrapping los bordes no existen.
        La grilla es un toroide.
  \item \textbf{Demostrar en vivo}: ejecutar \texttt{python juego\_vida.py}, cargar el
        blinker con la tecla 1, ver la oscilación.
  \item \textbf{El seed\_patterns.yml}: cómo las matrices binarias se convierten en
        conjuntos de coordenadas y se centran en la grilla.
\end{enumerate}

\newpage

% ================================================================
% SECCIÓN 2: REFINAMIENTO
% ================================================================
\section{Capa Refinamiento -- Miguel Ángel Torres Pérez}

\subsection{Archivos asignados}

\begin{itemize}
  \item \texttt{juego\_vida\_optimo.py} -- núcleo algorítmico puro, headless (203 líneas)
  \item \texttt{patterns.py} -- catálogo de 21 patrones + serialización YAML (259 líneas)
\end{itemize}

\subsection{Clase GameOfLifeOptimo}

Toma la lógica de \texttt{GameOfLife} y la refactoriza para que sea más limpia,
testeable y extensible. Las diferencias clave:

\begin{table}[h]
\centering
\begin{tabular}{lcc}
\toprule
\textbf{Aspecto} & \textbf{v1 (juego\_vida.py)} & \textbf{v2 (optimo.py)} \\
\midrule
Líneas de código         & 259 & 203 \\
Dependencias             & pygame, numpy, yaml & Solo yaml \\
GUI                     & Sí (800$\times$800) & No (headless) \\
Métodos puros           & 2 (parcial) & 4 (completos) \\
Manejo de patrones      & 9 patrones, código repetido & 21 patrones, API unificada \\
Reglas de autómata      & 1 (Conway) & 11 (todas) \\
API pública             & Solo process\_events/run\_logic/draw & 10 métodos programables \\
\bottomrule
\end{tabular}
\end{table}

\subsubsection{Métodos puros}

Un método puro es aquel que:
\begin{itemize}
  \item No modifica el estado del objeto (sin efectos secundarios)
  \item Solo depende de sus parámetros de entrada
  \item Es \textbf{thread-safe}: puede ejecutarse en paralelo sin locks
\end{itemize}

\begin{lstlisting}
def _count_living_neighbors(self, cell: Cell) -> int:
    """MÉTODO PURO: Sin efectos secundarios, thread-safe."""
    return sum(1 for n in self._get_neighbors(cell)
               if n in self._living_cells)

def _should_cell_survive(self, n: int) -> bool:
    """MÉTODO PURO: Encapsula reglas de supervivencia."""
    return self._rule.should_survive(n)

def _should_cell_reproduce(self, n: int) -> bool:
    """MÉTODO PURO: Encapsula reglas de reproducción."""
    return self._rule.should_reproduce(n)
\end{lstlisting}

\subsubsection{API programable}

A diferencia de la versión base (que solo se controla con teclas y mouse),
\texttt{GameOfLifeOptimo} expone una API que permite controlar todo desde código:

\begin{table}[h]
\centering
\begin{tabular}{ll}
\toprule
\textbf{Método} & \textbf{Uso} \\
\midrule
\texttt{load\_pattern("glider")}  & Carga y centra un patrón del catálogo \\
\texttt{run\_logic\_step()}       & Ejecuta 1 generación \\
\texttt{run\_n\_steps(500)}       & Ejecuta N generaciones \\
\texttt{set\_cells(cells)}        & Establece directamente las células vivas \\
\texttt{get\_cells()}             & Retorna el estado actual \\
\texttt{clear()}                  & Vacía la grilla \\
\texttt{set\_rule("highlife")}    & Cambia la regla de autómata \\
\texttt{list\_pattern\_names()}   & Lista los 21 patrones disponibles \\
\texttt{patterns\_by\_type("spaceship")} & Filtra patrones por categoría \\
\bottomrule
\end{tabular}
\end{table}

Ejemplo de uso:
\begin{lstlisting}
g = GameOfLifeOptimo()
g.load_pattern("r_pentomino")
g.run_n_steps(100)
print(f"Células vivas: {len(g.get_cells())}")
g.set_rule("highlife")         # cambiar regla en caliente
\end{lstlisting}

\subsection{\texttt{patterns.py} -- Catálogo de patrones}

\subsubsection{Clase Pattern}

\begin{lstlisting}
@dataclass
class Pattern:
    name: str              # "glider"
    type: str              # "spaceship"
    cells: Set[Cell]       # {(1,0), (2,1), (0,2), (1,2), (2,2)}
    period: int | None     # 4
    description: str       # "La nave más pequeña..."
    author: str            # "John Conway"
\end{lstlisting}

Propiedades y métodos clave:

\begin{table}[h]
\centering
\begin{tabular}{ll}
\toprule
\textbf{Método / Propiedad} & \textbf{Qué hace} \\
\midrule
\texttt{width}, \texttt{height} & Tamaño del bounding box \\
\texttt{bounds}                 & (min\_col, min\_row, max\_col, max\_row) \\
\texttt{to\_matrix()}           & Convierte Set[Cell] $\to$ matriz binaria normalizada \\
\texttt{centered(n\_cols, n\_rows)} & Centra el patrón en una grilla \\
\bottomrule
\end{tabular}
\end{table}

\subsubsection{Catálogo completo (21 patrones)}

\begin{table}[h]
\centering
\small
\begin{tabular}{llp{6cm}}
\toprule
\textbf{Categoría} & \textbf{Cant.} & \textbf{Patrones} \\
\midrule
Still life    & 7 & block, beehive, loaf, boat, tub, ship, pond \\
Oscilador     & 6 & blinker, toad, beacon, pulsar, pentadecathlon, queen\_bee \\
Nave espacial & 4 & glider, lwss, mwss, hwss \\
Cañón         & 1 & gosper\_glider\_gun (período 30, 36 células) \\
Matusalén     & 4 & r\_pentomino, diehard, acorn, b\_heptomino \\
\bottomrule
\end{tabular}
\end{table}

\subsubsection{Funciones de serialización YAML}

Convierten entre las tres representaciones del proyecto:

\begin{verbatim}
    Set[Cell]  ←── cells_to_matrix() / matrix_to_cells() ──→  List[List[int]]
         │                                                         │
         └──── save_pattern_to_yaml(cells, name, path) ────→  saved_patterns.yml
         └──── load_patterns_from_yaml(path) ←────────────  saved_patterns.yml
\end{verbatim}

\begin{lstlisting}
def cells_to_matrix(cells: Set[Cell]) -> List[List[int]]:
    """Convierte conjunto disperso a matriz binaria normalizada."""
    min_c = min(c for c,_ in cells)
    min_r = min(r for _,r in cells)
    w = max(c for c,_ in cells) - min_c + 1
    h = max(r for _,r in cells) - min_r + 1
    matrix = [[0] * w for _ in range(h)]
    for c, r in cells:
        matrix[r - min_r][c - min_c] = 1
    return matrix

def save_pattern_to_yaml(cells, name, filepath):
    """Guarda un patrón a YAML con ID autoincremental."""
    matrix = cells_to_matrix(cells)
    data = load_existing_yaml(filepath)
    next_id = max(data['patterns'].keys()) + 1
    data['patterns'][next_id] = {'name': name, 'matrix': matrix}
    yaml.dump(data, open(filepath, 'w'))
\end{lstlisting}

\subsection{Métricas de mejora}

\begin{itemize}
  \item \textbf{-31\% líneas} de código (259 $\to$ 203)
  \item \textbf{-75\% dependencias} externas (3 $\to$ 1)
  \item \textbf{+100\% métodos puros} (2 $\to$ 4)
  \item \textbf{+144\% patrones} soportados (9 $\to$ 21)
  \item \textbf{+1000\% reglas} de autómata (1 $\to$ 11)
\end{itemize}

\subsection{Puntos para exponer}

\begin{enumerate}
  \item \textbf{Métodos puros}: qué son y por qué son el primer paso hacia
        el paralelismo. Un método sin efectos secundarios puede ejecutarse
        en cualquier orden, en cualquier hilo.
  \item \textbf{API programable}: la versión base solo se usa con teclas.
        La versión refinada se puede controlar desde scripts, tests y otras
        aplicaciones.
  \item \textbf{El catálogo de patrones}: 21 patrones con metadatos.
        Cada patrón sabe su tamaño, su período, su autor. Demo: cargar
        un glider, un pulsar, un gosper\_glider\_gun.
  \item \textbf{Serialización YAML}: cómo un patrón dibujado con el mouse
        puede guardarse a disco y cargarse después. El flujo completo:
        Set[Cell] $\to$ matriz $\to$ YAML $\to$ disco.
  \item \textbf{Métricas}: mostrar la tabla de mejora. Menos código,
        más funcionalidad.
\end{enumerate}

\newpage

% ================================================================
% SECCIÓN 3: ESCALABILIDAD
% ================================================================
\section{Capa Escalabilidad -- Andrés Soria Cabrera}

\subsection{Archivos asignados}

\begin{itemize}
  \item \texttt{juego\_vida\_proyecto.py} -- versión paralela con multiprocessing (210 líneas)
  \item \texttt{automata.py} -- motor de 11 reglas de autómata configurables (65 líneas)
\end{itemize}

\subsection{\texttt{juego\_vida\_proyecto.py} -- Motor paralelo}

\subsubsection{¿Por qué funciones sueltas y no una clase?}

Python requiere que las funciones enviadas a otros procesos vía \texttt{multiprocessing}
sean \textbf{serializables con pickle}. Las funciones definidas a nivel de módulo
(no como métodos de clase) cumplen este requisito. Por eso \texttt{calcular\_segmento}
y \texttt{contar\_vecinos} son funciones globales.

\subsubsection{Conversión Set $\leftrightarrow$ Grid}

La versión paralela usa una matriz densa \texttt{List[List[int]]} de 200$\times$200
en vez de \texttt{Set[Cell]}. Las funciones de conversión son:

\begin{lstlisting}
def cells_set_to_grid(cells, filas, cols, offset_col=0, offset_row=0):
    """Set[Cell] -> matriz 2D con offset opcional."""
    tablero = [[0]*cols for _ in range(filas)]
    for c, r in cells:
        if 0 <= r+offset_row < filas and 0 <= c+offset_col < cols:
            tablero[r+offset_row][c+offset_col] = 1
    return tablero

def grid_to_cells_set(tablero):
    """Matriz 2D -> Set[Cell] (para guardar a YAML)."""
    cells = set()
    for r, row in enumerate(tablero):
        for c, val in enumerate(row):
            if val: cells.add((c, r))
    return cells
\end{lstlisting}

\subsubsection{El worker: \texttt{calcular\_segmento()}}

Esta función es el núcleo del paralelismo. Cada proceso trabajador recibe
una \textbf{franja de filas} y calcula la siguiente generación para esa franja.

\begin{lstlisting}
def calcular_segmento(args):
    """Worker de multiprocessing. Calcula una franja de filas."""
    tablero, fila_inicio, fila_fin, filas_tot, cols_tot, born, survive = args
    segmento_nuevo = []

    for fila in range(fila_inicio, fila_fin):
        nueva_fila = []
        for col in range(cols_tot):
            vecinos = contar_vecinos(tablero, col, fila, filas_tot, cols_tot)
            estado_actual = tablero[fila][col]

            # Reglas de Conway con born/survive dinámicos
            if estado_actual == 1 and vecinos not in survive:
                nueva_fila.append(0)      # muere
            elif estado_actual == 0 and vecinos in born:
                nueva_fila.append(1)      # nace
            else:
                nueva_fila.append(estado_actual)  # no cambia

        segmento_nuevo.append(nueva_fila)

    return fila_inicio, fila_fin, segmento_nuevo
\end{lstlisting}

\textbf{Datos clave del worker:}
\begin{itemize}
  \item Recibe \texttt{born} y \texttt{survive} como sets de enteros $\to$ reglas dinámicas
  \item Solo lee el tablero actual (nunca lo modifica) $\to$ sin condiciones de carrera
  \item Escribe exclusivamente en su propio \texttt{segmento\_nuevo} $\to$ sin conflicto
  \item Usa \texttt{contar\_vecinos()} con wrapping toroidal siempre activo
\end{itemize}

\subsubsection{Pool.map() y descomposición por filas}

El trabajo se divide en franjas horizontales, una por núcleo de CPU:

\begin{lstlisting}
num_procesos = os.cpu_count() - 1     # Ej: 7 procesos en CPU de 8 núcleos
pool = mp.Pool(processes=num_procesos)

filas_por_proc = 200 // num_procesos  # Ej: 200/7 = 28 filas por proceso
tareas = []
for i in range(num_procesos):
    inicio = i * filas_por_proc
    fin = 200 if i == num_procesos-1 else (i+1) * filas_por_proc
    tareas.append((tablero, inicio, fin, 200, 200, rule.born, rule.survive))

# EJECUCIÓN PARALELA
resultados = pool.map(calcular_segmento, tareas)

# REENSAMBLADO
nuevo_tablero = crear_tablero_vacio()
for inicio, fin, segmento in resultados:
    nuevo_tablero[inicio:fin] = segmento
tablero = nuevo_tablero
\end{lstlisting}

\begin{center}
\begin{tikzpicture}[scale=0.7,
  strip/.style={rectangle,draw,fill=blue!10,text width=2cm,text centered,rounded corners,font=\tiny}]
  \draw (0,0) rectangle (6,4);
  \node at (3,4.3) {Tablero 200$\times$200};

  \node[strip,fill=red!15] at (3,0.5) {CPU 0: filas 0--49};
  \node[strip,fill=green!15] at (3,1.5) {CPU 1: filas 50--99};
  \node[strip,fill=yellow!15] at (3,2.5) {CPU 2: filas 100--149};
  \node[strip,fill=blue!15] at (3,3.5) {CPU 3: filas 150--199};

  \draw[->,>=latex,thick] (6.5,2) -- (7.5,2) node[midway,above] {\tiny Pool.map};
  \draw (7.5,0) rectangle (9.5,4);
  \node at (8.5,4.3) {Nuevo tablero};
\end{tikzpicture}
\end{center}

\subsubsection{Cámara y renderizado}

\begin{itemize}
  \item \textbf{Zoom}: rueda del mouse, 4px -- 60px por célula
  \item \textbf{Paneo}: teclas WASD o flechas
  \item \textbf{Renderizado selectivo}: solo se dibujan las celdas visibles
  \item \textbf{Grilla}: se oculta con zoom $\leq$ 5px por rendimiento
  \item \textbf{Pintado}: click izq = vivir, click der = matar (solo en pausa)
  \item Las coordenadas del mouse se ajustan por cámara y zoom
\end{itemize}

\subsection{\texttt{automata.py} -- Motor de reglas}

\subsubsection{Clase AutomataRule}

\begin{lstlisting}
@dataclass
class AutomataRule:
    born: Set[int]       # Ej: {3} para Conway
    survive: Set[int]    # Ej: {2, 3} para Conway
    name: str            # "Conway's Life"

    def should_survive(self, n: int) -> bool:
        return n in self.survive

    def should_reproduce(self, n: int) -> bool:
        return n in self.born

    @property
    def rule_string(self) -> str:
        return f"B{born}/S{survive}"   # "B3/S23"
\end{lstlisting}

\subsubsection{11 presets}

\begin{table}[h]
\centering
\small
\begin{tabular}{llp{6cm}}
\toprule
\textbf{Preset} & \textbf{Regla} & \textbf{Comportamiento} \\
\midrule
conway     & B3/S23      & Clásico de Conway: gliders, osciladores, cañones \\
highlife   & B36/S23     & Parecido a Conway pero con replicador natural \\
seeds      & B2/S        & Todo es caótico, los patrones mueren rápido \\
maze       & B3/S12345   & Genera estructuras tipo laberinto \\
daynight   & B3678/S34678& Simétrico: invertir vivo/muerto da igual \\
replicator & B1357/S1357 & Patrones que se copian a sí mismos \\
morley     & B368/S245   & Naves de alta velocidad y explosiones \\
diamoeba   & B35678/S5678& Formas ameboides que se expanden \\
2x2        & B36/S125    & Dominado por bloques 2$\times$2 \\
34life     & B34/S34     & Explosivo e impredecible \\
longlife   & B345/S5     & Estructuras que persisten mucho tiempo \\
\bottomrule
\end{tabular}
\end{table}

\subsubsection{Parser de notación B/S}

\begin{lstlisting}
def parse_rule(rule_string: str) -> AutomataRule:
    """'B3/S23' -> AutomataRule(born={3}, survive={2,3})"""
    parts = rule_string.upper().split("/")
    born = set()
    survive = set()
    for part in parts:
        if part.startswith("B"): born = {int(d) for d in part[1:]}
        if part.startswith("S"): survive = {int(d) for d in part[1:]}
    return AutomataRule(born, survive)

def get_rule(name_or_string: str) -> AutomataRule:
    """Busca en PRESETS; si no, parsea como regla personalizada."""
    if name_or_string in PRESETS:
        return PRESETS[name_or_string]
    return parse_rule(name_or_string)
\end{lstlisting}

\subsection{Puntos para exponer}

\begin{enumerate}
  \item \textbf{Paralelismo de datos}: el tablero se divide en franjas de filas.
        Cada CPU calcula su franja sin comunicarse con las demás. Al final se reensambla.
        Es el mismo algoritmo, pero distribuido.
  \item \textbf{¿Por qué funciona?}: cada generación depende solo de la generación
        anterior. No hay dependencia entre franjas durante el cálculo. Cada worker
        solo lee el tablero completo y escribe en su propio segmento.
  \item \textbf{Las reglas no son fijas}: el worker recibe \texttt{born} y \texttt{survive}
        como parámetros. Cambiar de Conway a HighLife es instantáneo.
  \item \textbf{11 universos distintos}: mismo patrón, diferentes reglas = comportamientos
        radicalmente distintos. Demo: glider con Conway $\to$ se mueve; con Seeds $\to$ muere;
        con Day\&Night $\to$ algo completamente diferente.
  \item \textbf{Cámara y zoom}: la grilla real es de 200$\times$200, pero la pantalla
        es de 800$\times$800. La cámara permite explorar el tablero virtual.
\end{enumerate}

\newpage

% ================================================================
% APÉNDICE
% ================================================================
\section{Apéndice: referencia rápida de comandos}

\subsection{Ejecutar cada versión}

\begin{verbatim}
# Versión base (GUI)
python juego_vida.py

# Versión optimizada (demo automática)
python juego_vida_optimo.py

# Versión paralela (GUI con cámara)
python juego_vida_proyecto.py
\end{verbatim}

\subsection{Cargar patrones desde código}

\begin{verbatim}
from juego_vida_optimo import GameOfLifeOptimo
g = GameOfLifeOptimo()

# Patrones disponibles
g.list_pattern_names()

# Cargar y ejecutar
g.load_pattern("glider")
g.run_n_steps(10)
print(g.get_cells())

# Cambiar regla
g.set_rule("highlife")
g.load_pattern("r_pentomino")
g.run_n_steps(500)
\end{verbatim}

\subsection{Controles unificados (versiones GUI)}

\begin{table}[h]
\centering
\begin{tabular}{ll}
\toprule
\textbf{Tecla} & \textbf{Acción} \\
\midrule
Espacio   & Pausar / reanudar \\
C         & Limpiar grilla \\
1--9      & Cargar patrón precargado \\
P         & Imprimir catálogo en consola \\
R         & Cambiar regla de autómata \\
S         & Guardar patrón actual a YAML \\
Click     & Pintar / borrar célula \\
Rueda     & Zoom (solo v3) \\
WASD      & Mover cámara (solo v3) \\
\bottomrule
\end{tabular}
\end{table}

\end{document}
